% Invented text. A representative article: a title block and abstract,
% sections of justified prose that hyphenates, footnotes, nested lists, a
% booktabs table, a figure, display and inline mathematics, a quotation, a
% description list, and an appendix with a code block whose longest line
% does not fit and a long URL — the mixed pages an article sets.
\documentclass{article}
\usepackage{graphicx}
\usepackage{booktabs}
\usepackage{hyperref}
\fonts{ dir = "../corpus/fonts", body = "Source Serif Pro", sans = "Fira Sans",
  mono = "Source Code Pro", math = "Fira Math" }

\title{Rainfall on a Terraced Garden}
\author{A. Placeholder}
\date{}

\begin{document}

\maketitle

\begin{abstract}
Four seasons of rain gauge readings from an invented hillside garden are
compared against the soil moisture recorded on each of its three
terraces. The upper terrace drains within hours, the lower one holds water
for days, and the middle terrace behaves like neither: its response depends
on how much rain fell during the preceding week. A simple storage model
with one parameter per terrace reproduces the readings closely enough to
plan irrigation from it.
\end{abstract}

\section{Introduction}

A garden built on a slope is three gardens. Water that falls on the upper
terrace runs through it and arrives, later and slower, at the terraces
below; what each bed receives is therefore not the rainfall but a
transformation of it, smoothed and delayed by every bed above. Gardeners
know this without measuring it, and plant accordingly: the thirsty crops
go to the bottom, the herbs that dislike wet roots go to the top. The
question this note asks is narrower. Given a season of rainfall, how much
water reaches each terrace, and when?\footnote{The garden, its terraces
and every number below are invented for this document.}

The measurements were made with the simplest instruments available: a
funnel gauge emptied each morning, and a resistive moisture probe pushed
into each terrace at a fixed depth. Neither instrument is precise, but
both are consistent, and consistency is what a comparison between
terraces requires.

\subsection{The site}

The garden faces south-east and drops about four metres over its length.
Its three terraces are retained by dry-stone walls, which leak freely; no
terrace is sealed from the one below it. The soil is a sandy loam on the
upper terrace, grading into a heavier clay loam at the bottom, where
decades of runoff have carried the fine particles down.

\begin{itemize}
  \item The upper terrace holds the gauge, a row of rosemary and the
  compost bins.
  \item The middle terrace is the vegetable garden proper:
  \begin{itemize}
    \item beans and peas along the wall;
    \item brassicas in the centre beds, rotated yearly;
    \item a strip of salad leaves at the shaded end.
  \end{itemize}
  \item The lower terrace is lawn, with a pond in its wettest corner.
\end{itemize}

\section{Method}

Readings were taken every morning at the same hour for four seasons. A
missed morning was recorded as missing rather than interpolated, and the
model below was fitted only to complete weeks.

\begin{enumerate}
  \item Empty the gauge and record the depth of water in millimetres.
  \item Read the probe on each terrace, upper first, wiping the
  electrodes between readings.
  \item Note anything unusual: frost on the probe, a blocked funnel,
  standing water on the lawn.
\end{enumerate}

Each terrace is modelled as a store that fills with the water arriving
from above and empties in proportion to what it holds. With $S_t$ the
water stored on day $t$, $R_t$ the inflow and $k$ the terrace's drainage
fraction,
\begin{equation}
  S_{t+1} = (1 - k)\, S_t + R_t ,
  \qquad 0 < k < 1 ,
\end{equation}
and the outflow $k S_t$ is the inflow of the terrace below. The upper
terrace receives the rainfall itself; the others receive what the terrace
above releases.

\section{Results}

Table~\ref{tab:fit} gives the fitted drainage fraction of each terrace,
with the share of the variation in its moisture readings that the model
explains. The upper terrace drains fastest, as expected, and the model
fits it best: its store is small and its response immediate.

\begin{table}
\centering
\begin{tabular}{@{}lrrp{0.36\linewidth}@{}}
  \toprule
  \textbf{Terrace} & \textbf{$k$} & \textbf{Fit} & \textbf{Remark} \\
  \midrule
  Upper & 0.62 & 0.91 & drains within a day of heavy rain \\
  Middle & 0.31 & 0.78 & depends strongly on the preceding week \\
  Lower & 0.12 & 0.84 & stays wet through most of the winter \\
  \bottomrule
\end{tabular}
\caption{Drainage fraction and fit for each terrace.}
\label{tab:fit}
\end{table}

The middle terrace is the interesting case. In a dry week it behaves like
the upper terrace, taking water in and passing it on quickly; after a wet
week it behaves like the lower one, holding what it receives. A single
fraction cannot describe both, and its fit is correspondingly worse. The
residuals show the pattern clearly: the model underestimates the middle
terrace after every run of wet days and overestimates it after every dry
spell.

\begin{figure}
\centering
\includegraphics[width=0.5\textwidth, alt={Three coloured rectangles
standing for the three terraces}]{../corpus/rects.png}
\caption{The three terraces, schematically: upper, middle and lower.}
\end{figure}

A two-store model for the middle terrace, one store for the topsoil and
one for the clay beneath it, removes most of this structure. Its extra
parameter is the fraction $f$ of the inflow that reaches the deeper store
directly, and with $f \approx 0.4$ the fit rises to $0.89$, close to the
other terraces.

\section{Discussion}

The models are crude, and the instruments cruder, but the comparison they
support is robust: the terraces differ by a factor of five in how quickly
they drain, and the middle one remembers the weather. An old gardening
manual puts the same observation more briefly:

\begin{quote}
Water the top of a slope often and the bottom seldom; the middle will
tell you which it wants, if you watch it for a week.
\end{quote}

\begin{description}
  \item[Irrigation] On the evidence here, the upper terrace needs water
  within two dry days in summer, the lower terrace hardly ever.
  \item[Planting] Crops that tolerate both drought and saturation belong
  on the middle terrace, or nowhere.
  \item[Measurement] A second probe at a greater depth on the middle
  terrace would test the two-store model directly.\footnote{The probe
  has been ordered; its readings will be compared in a later season.}
\end{description}

\appendix

\section{Reading the logs}

The daily readings are kept as plain text, one line per morning. The
script that fits the drainage fractions reads them as follows; its longest
line is longer than this measure, and is kept as written.

\begin{verbatim}
def read_log(path):
    rows = []
    for line in open(path):
        day, rain_mm, upper, middle, lower = line.split(",")
        rows.append((int(day), float(rain_mm), float(upper), float(middle), float(lower)))
    return rows
\end{verbatim}

The logs and the fitting script are kept with the garden's other records at
\url{https://example.org/gardens/terraced-hillside/records/rainfall-and-soil-moisture/2026/daily-log.txt},
where each season's file is added when the season ends.

\end{document}
